\chapter{从原子到星系：宇宙的尺度}\label{ch:g9:atoms-to-galaxies}

九年前，这本书从你的五官和一只瓢虫的腿开始。
它的结尾是一段单凭感官无法完成的旅程：
跨越四十一个数量级，从原子的核心
直到可观测宇宙的边缘——而你，确切地说，
正站在正中间。最后一章，兑现一个老笑话欠下的秘密，
并搭起一架梯子，
供此后的全部物理学攀登。

\section{向下：深入物质}

\begin{definition}[原子]\label{def:g9:atoms-to-galaxies:atom}
你七岁时接触的分子本身也有结构：
每个分子都由\emph{原子}\index{原子}组成——这是自然界的字母表，
约一百种，它们的组合拼写出每一种物质。
水分子：两个氢原子与一个氧原子。
原子本身又有结构：
中心是一个极小、很重、带正电的
\emph{原子核}\index{原子核}，远离中心的外围是一团
\emph{电子}\index{电子}——
几乎没有重量，是带负电的载体。原子的大小约为
$10^{-10}$ \unit{m}，原子核的大小约为 $10^{-15}$ \unit{m}，
又小了约十万倍。
\end{definition}

\begin{proposition}[物质主要是空的]\label{prop:g9:atoms-to-galaxies:empty}
把一个原子放大到一座大体育场那么大，原子核是中圈中央的一粒豌豆，
电子则是最高看台上的几只蚊子：其间全部是空的。
大理石桌、铁砧、你自己的手，绝大部分都是空的，
它们的坚实感来自一种电的“排斥”，
原子的电子云相互排斥，造成充实的假象。
这是物理学中最令人惊讶的句子之一，
其内容不过是
$10^{-10}$ 与 $10^{-15}$ 的简单算术。
\end{proposition}

\begin{example}[旧笑话的兑现]\label{ex:g9:atoms-to-galaxies:electron}
电学章节曾欠你一个小笑话，现在兑现。
本书每一根导线中运动的电荷，
正是\emph{电子}——它们带负电，
因此从电池的 $-$ 接线柱走向 $+$ 接线柱，
恰好与一个世纪前物理学家规定的电流方向相反，
那时还没有人能询问这些运动的电荷。
这一约定被保留了下来——你电路中的每一条规律
在这一约定下都完美成立——但这小小的讽刺依然存在：在你画过的每一张电路图中，
真正的“人群”走的是另一个方向。
\end{example}

\begin{example}[再往下]\label{ex:g9:atoms-to-galaxies:deeper}
原子核还有自己的组成部分，它们如何结合，
为何有些原子核会裂变（发电站的热量来源），
或聚变（太阳的能量来源）——这些故事把物理学一直带到 $10^{-15}$
米甚至更小，进入人类所探测到的最微小结构。
高中分册会正式展开原子核的内容；
大学阶段则会走得更远。这架梯子的较低横档
仍在建造之中。
\end{example}

\section{向上：深入天空}

\begin{example}[攀登]\label{ex:g9:atoms-to-galaxies:ascent}
现在开始攀登，一个数量级接一个数量级，经过
九年来熟悉的各个地标。你，约为 $10^{0}$ \unit{m}。
学校操场，$10^{2}$。地球的直径为 $10^{7}$ \unit{m}——
正是体育场模型中的那颗弹珠。月球的距离为
$4 \times 10^{8}$ \unit{m}；太阳的距离为 $1.5 \times 10^{11}$ \unit{m}——约八光分。
太阳系的宽度约为 $10^{13}$ \unit{m}；
最近的恒星距离为 $4 \times 10^{16}$ \unit{m}——约四光年。
银河系的大小约为 $10^{21}$ \unit{m}，即一座城市的尺度；
邻近的星系距离为 $2 \times 10^{22}$ \unit{m}；
我们能观测到的最深天空约为 $10^{26}$ \unit{m}——
那里星系发出的光，比你椅子底下的地球还要古老。
\end{example}

\begin{omfigure}
\begin{tikzpicture}[scale=0.85]
  \draw[-Stealth, very thick] (0,0.6) -- (11.4,0.6);
  \foreach \x/\p in {0.4/-15, 1.7/-10, 3.0/-5, 4.3/0, 5.6/5,
      6.9/10, 8.2/15, 9.5/20, 10.8/25} {
    \draw[thick] (\x,0.45) -- (\x,0.75);
    \node[below, font=\footnotesize] at (\x,0.4) {$10^{\p}$};
  }
  % labels on two alternating rows, with thin leaders
  \foreach \x/\h/\name in {0.4/0.95/原子核, 1.7/1.6/原子,
      3.0/0.95/蚂蚁, 4.3/1.6/你, 5.6/0.95/{一个国家},
      6.9/1.6/{日地距离}, 8.2/0.95/{最近的恒星},
      9.5/1.6/{银河系}, 10.8/0.95/{深空}} {
    \draw[omExo!80, thin] (\x,0.75) -- (\x,\h);
    \node[above, font=\footnotesize] at (\x,\h) {\name};
  }
\end{tikzpicture}

{\small 尺度阶梯（单位：米）：从原子核到深空共四十一个
数量级——
读者几乎恰好站在正中间的横档上。（蚂蚁原谅了我们的
取整。）}
\end{omfigure}

\begin{method}[用数量级思考]\label{met:g9:atoms-to-galaxies:oom}
物理学家面对任何新问题时的第一个工具：
\begin{enumerate}
  \item 用科学记数法表示该量；
  \item 只保留 10 的幂次——\emph{数量级
        }\index{数量级}；
  \item 比较梯级而非数字：位于 $10^{16}$
        米的恒星与位于 $10^{11}$ 米的行星相差五个横档，
        ——即相差十万倍，
        无论前导数字如何；
  \item 只有在需要时，再细化前导数字。
\end{enumerate}
九年来每一步的判断——“这个答案合理吗？”——
——逐渐凝成了这一习惯。高中分册的第一章
正是在此基础上建立它的测量思维。
\end{method}

\begin{example}[梯级运算]\label{ex:g9:atoms-to-galaxies:arithmetic}
多少个原子能排满一支铅笔的宽度？铅笔宽 $10^{-2}$
\unit{m}；原子 $10^{-10}$：相差八个梯级——
一亿个原子，肩并肩。从地球到太阳，能排下多少个地球？
$1.5 \times 10^{11}$ 除以 $1.3 \times 10^{7}$，约为
$10^{4}$——即运动场模型中的一万颗弹珠，
排成一行。梯级法能在几秒内给出答案，而用原始的
千米数则被一连串的零淹没。
\end{example}

\section{自中间梯级望去}

\begin{remark}[九年所建]\label{rem:g9:atoms-to-galaxies:built}
回头看看这条梯子，数一数你所拥有的。在
$10^{-10}$ 处：分子与原子，连同它们带来的三态、
热量、声音的接力、电流中的电子。约 $10^{0}$ 处：
力与能量、光的直线传播、电路的规律——
这是九年实验所得的、与人体尺度相当的物理。从
$10^{7}$ 到 $10^{22}$：自转的地球、倾斜的季节、
下落的月球，以及一条同时支配果园与银河的万有引力定律。
没有哪一门学科，能让你从物质内部，
一路领你到可见世界的边缘——
而每一级都借助仪器、诚实，
以及你可以亲自核验的算术，
攀登而得。
\end{remark}

\begin{remark}[前路所待]\label{rem:g9:atoms-to-galaxies:ahead}
未竟之事，恰是其中最精彩的部分。究竟是什么定律
把力与运动紧紧联系在一起？光到底\emph{是}什么，
它竟能以 $3 \times 10^{8}$ 米每秒飞奔，并在其中携带颜色？
原子核内部发生着什么样的裂变——
又是什么锻造了太阳光？为什么有用能量的账目
总是向下滑落？每一个问题，本书都曾诚实地提出，
也诚实地搁置；每个问题都有答案，
而这些答案正是人类所知的最精妙的内容之一。
它们将在一年后的高中分册中展开。
带上这把梯子，
也带上凡事都要核验的习惯。
\end{remark}

\section{练习}

\begin{exercise}[$\star$]\label{exo:g9:atoms-to-galaxies:1}
按大小排序：分子、原子、原子核、蚂蚁。
分别给出各自以米为单位的 10 的幂次。
\end{exercise}

\begin{exercise}[$\star$]\label{exo:g9:atoms-to-galaxies:2}
描述原子的结构：位于中心的是什么，
围绕它的又是什么，以及构成体育场模型的
两个尺度。
\end{exercise}

\begin{exercise}[$\star$]\label{exo:g9:atoms-to-galaxies:3}
把那个老笑话继续讲下去：导线中究竟是谁在行进，
朝哪个方向——揭示之后，为什么电路的一切规律
依然成立？
\end{exercise}

\begin{exercise}[$\star$]\label{exo:g9:atoms-to-galaxies:4}
什么是数量级？给出下列各量的数量级：
你的身高、教室的长度、
地球的直径。
\end{exercise}

\begin{exercise}[$\star$]\label{exo:g9:atoms-to-galaxies:5}
这些量相差几个梯级：原子与原子核，
你与地球，太阳的距离与
最近恒星的距离？
\end{exercise}

\begin{exercise}[$\star$]\label{exo:g9:atoms-to-galaxies:6}
“桌子大部分是空的。”请借助两个原子尺度的数据
为这句话辩护——并解释桌子为何仍然
让人感觉坚实。
\end{exercise}

\begin{exercise}[$\star$]\label{exo:g9:atoms-to-galaxies:7}
梯级运算：大约多少个原子排成一行，
才能铺满一毫米？
\end{exercise}

\begin{exercise}[$\star$]\label{exo:g9:atoms-to-galaxies:8}
在梯子上定位以下各量，并各给出一个 10 的幂次：光年、
银河、已观测的最深天空。
\end{exercise}

\begin{exercise}[$\star\star$]\label{exo:g9:atoms-to-galaxies:9}
体育场模型：若原子核是一颗 \qty{1}{cm} 的豌豆，
那么原子这座体育场有多宽？（比值为 $10^{5}$——换算为
米，并与真实的体育场作比较。）
\end{exercise}

\begin{exercise}[$\star\star$]\label{exo:g9:atoms-to-galaxies:10}
一滴水约含 $10^{21}$ 个分子。利用七年级那一章中的
“水滴与海洋”比较——
或者你自己对全部海洋中水滴数的估计（约
$10^{21}$ 升，每升含 $10^{5}$ 滴）——权衡一下这条
古老说法：一滴水所含的分子，比大海中所含的
水滴还多。
\end{exercise}

\begin{exercise}[$\star\star$]\label{exo:g9:atoms-to-galaxies:11}
你的身体位于梯子的中部附近：向下到原子、
向上到最深天空，大约各有几级？这种对称性
说明了什么，关于九年学业所能抵达的
范围？
\end{exercise}

\begin{exercise}[$\star\star\star$]\label{exo:g9:atoms-to-galaxies:12}
这是本书的最后一道练习。从你九年的学习中任选三章——
一章来自一至三年级，一章来自四至六年级，
一章来自七至九年级——对每一章，写下一句话，
告诉开始这一章时的年轻的自己，
从梯子的顶端回望，它真正教会了你什么。
（写给自己的信没有参考答案——
但请动笔写吧；接下来的课程会为你感到高兴。）
\end{exercise}

\section{问题：十的幂次之旅}

\begin{problem}[{周末问题 --- 班级拍摄“十的幂次
之旅”一片；一次连续变焦，跨越四十一级；
剧本研讨会}]\label{pb:g9:atoms-to-galaxies:1}
班级的影片：一次连续变焦，从野餐毯
向外推至深空，再向内直至原子核——
影片每秒对应一个十的幂次。你来撰写
科学剧本。

\medskip\noindent\textbf{第 I 部分 --- 向外。}
变焦开始时，画面框住一米长的野餐毯：$10^{0}$。
\begin{enumerate}
  \item 画面在哪些秒数首次包含：
        整个校园（$10^{2}$）；整个地球
        （$10^{7}$）；月球轨道（$10^{9}$）？
  \item 哪一秒画面框住了太阳的距离——
        依据讲光的那一章，旁白说阳光
        的年龄是多少？
  \item 画面在哪两秒之间穿过那片
        巨大的空虚——太阳系已经到头，
        而最近的恒星还未进入？屏幕上
        应该如实呈现什么？
  \item 变焦在第 $26$ 秒结束：画面中充满了什么，
        其中最古老的光与地球相比，
        有多古老？
\end{enumerate}

\medskip\noindent\textbf{第 II 部分 --- 向内。}
变焦方向反转，潜入毯子上的一片树叶。
\begin{enumerate}[resume]
  \item 影片在哪几秒经过：一只蚂蚁
        （$10^{-2}$，宽泛估计）；一个细胞（$10^{-5}$）；
        一个分子（$10^{-9}$）；一个原子（$10^{-10}$）？
  \item 在原子的电子云与原子核之间，
        影片还必须再变焦五秒穿过——
        穿过的是什么？旁白对
        “体育场”怎么说？
  \item 影片的最后一幅画面位于 $10^{-15}$：
        其主体是什么——以及旁白关于
        下方各级的坦率结语。
  \item 计算影片总长：向外几秒，
        向内几秒，从毯子的
        一米算起？
\end{enumerate}

\medskip\noindent\textbf{第 III 部分 --- 剧本
研讨会。}
\begin{enumerate}[resume]
  \item 一位老师异议：影片在从房屋到街道
        的一步上用了一秒，在从星系到星系团
        的一步上也用了一秒——“步长极不均匀！”
        请用梯级的逻辑
        为等长的秒数辩护。
  \item 旁白希望两个方向共用一句反复出现的台词——
        在每一级上关于空虚都成立的一句话。
        请起草这句话，
        它须同时适用于太阳系与原子。
  \item 片尾字幕须列出征服两个方向的两件
        仪器，请写出它们的名称，
        并指出本书中
        介绍其原理的章节。
  \item 影片的结束语属于课程本身：
        请写两句话的尾声——九年，
        一架梯子，以及明年即将开始的内容。
\end{enumerate}
\end{problem}
